% !TEX program = xelatex
% -*- coding: utf-8 -*-
% 状元王加赛纸：独立一张，只服务决赛桌；各桌规则纸不含状元王内容。
% 状元等级表与各桌规则纸第三节逐字一致，由 scripts/version_consistency.py 跨件比对。
\documentclass[11pt,a4paper,fontset=none]{ctexart}
\usepackage[left=1.65cm,right=1.65cm,top=1.25cm,bottom=1.25cm]{geometry}
\usepackage{tikz}
\usepackage{array,tabularx}
\usepackage[table]{xcolor}
\usepackage{enumitem}
\usepackage{qrcode}
\setCJKmainfont{Songti SC}
\setCJKsansfont{Hiragino Sans GB}
\setCJKmonofont{Hiragino Sans GB}
\pagestyle{empty}
\linespread{1.06}
\setlength{\parindent}{0pt}
\setlength{\parskip}{0pt}
\setlength{\tabcolsep}{6pt}
\setlist[enumerate]{leftmargin=1.6em,label=\textbf{\arabic*.},itemsep=2pt,topsep=3pt,parsep=0pt}
\definecolor{headgray}{gray}{0.16}
\definecolor{rowgray}{gray}{0.95}
\definecolor{rulegray}{gray}{0.72}
\newcolumntype{L}{>{\raggedright\arraybackslash}X}
\newcommand{\shead}[2]{\par\vspace{7pt}%
  {\fontsize{12.5}{17}\selectfont\bfseries\sffamily
  #1、#2\hspace{10pt}{\color{rulegray}\leaders\hrule height 0.3pt\hfill}\kern0pt\par}\vspace{4pt}}
\newcommand{\tablehead}[1]{\textcolor{white}{\textsf{\textbf{#1}}}}
\newcommand{\die}[1]{%
  \begin{tikzpicture}[baseline=0.5ex]
    \ifnum#1=4
      \filldraw[black,line width=0.5pt,rounded corners=0.8pt] (0,0) rectangle (0.40,0.40);
      \node[text=white,inner sep=0pt,font=\sffamily\bfseries\fontsize{9}{9}\selectfont] at (0.20,0.20) {#1};
    \else
      \draw[line width=0.5pt,rounded corners=0.8pt] (0,0) rectangle (0.40,0.40);
      \node[inner sep=0pt,font=\sffamily\fontsize{9}{9}\selectfont] at (0.20,0.20) {#1};
    \fi
  \end{tikzpicture}%
}
\newcommand{\roll}[6]{\die{#1}\,\die{#2}\,\die{#3}\,\die{#4}\,\die{#5}\,\die{#6}}

\begin{document}
\fontsize{11}{16}\selectfont
\begin{minipage}[c]{0.79\linewidth}
  {\fontsize{28}{34}\selectfont\bfseries 状元王加赛规则\par}
  \vspace{7pt}
  {\fontsize{10.5}{14.5}\selectfont 本次活动统一版\quad |\quad 各桌状元参加，决出状元王\par}
\end{minipage}\hfill
\begin{minipage}[c]{0.18\linewidth}
  \centering
  \qrcode[height=20mm,level=M,padding]{https://www.luochang.ink/bobing-game/champion-final/}\par
  {\fontsize{9}{11}\selectfont\sffamily 扫码看规则\par}
\end{minipage}\par
\vspace{5pt}\hrule height 0.8pt

\shead{一}{参加与座次}
\begin{enumerate}
  \item \textbf{参赛}：各桌状元同台围一桌，抽签定掷骰顺序。
  \item \textbf{轮次}：顺时针轮流，每人一掷，掷完传左手边；所有人各掷一次为一轮。
  \item \textbf{记录}：记录员由玩家兼任，记下当前领先者的六颗骰面。
\end{enumerate}

\shead{二}{判定与当选}
每一掷当场与全场最好成绩比较：\textbf{先比等级，再比带点}；等级与带点都相同，先掷出者保留领先。每人只记本场最好成绩；未掷出状元的掷骰不参与比较。黑底“4”表示四点，俗称“红四”。表中骰面是一组示例，排列顺序不限，判定以文字条件为准。

\vspace{3pt}
{\renewcommand{\arraystretch}{1.24}
\begin{tabularx}{\linewidth}{p{2.35cm} >{\centering\arraybackslash}p{3.1cm} L}
\rowcolor{headgray}\tablehead{状元等级} & \tablehead{骰面示例} & \tablehead{条件／同级比较}\\
状元插金花 & \roll{4}{4}{4}{4}{1}{1} & 4 颗四点＋2 颗一点；本场最高\\
\rowcolor{rowgray}六红 & \roll{4}{4}{4}{4}{4}{4} & 6 颗四点\\
六同 & \roll{2}{2}{2}{2}{2}{2} & 6 颗同点，四点除外，一点也算；彼此同级\\
\rowcolor{rowgray}五红 & \roll{4}{4}{4}{4}{4}{6} & 恰有 5 颗四点；比剩余 1 颗，大者胜\\
五子登科 & \roll{2}{2}{2}{2}{2}{6} & 恰有 5 颗同点，非 4；只比剩余 1 颗\\
\rowcolor{rowgray}四红 & \roll{4}{4}{4}{4}{2}{6} & 恰有 4 颗四点，排除插金花；比另 2 颗之和\\
\end{tabularx}\par}
\vspace{3pt}
\textbf{当选：}领先成绩产生后，从其下一位起，其他每位参赛者各掷一次而无人超越时，领先者当选状元王，加赛结束。领先者本人掷出更好的状元，重新起算。

\shead{三}{连续无状元时}
\textbf{连续 10 轮无人掷出任何状元}，则每人加掷一次：掷出状元的按第二节分出名次；都未掷出，比六颗点数总和，最大者当选；总和相同者再掷，直至分出。

\shead{四}{小提醒}
\textbf{出碗：}任何一颗未入碗或跳出碗，这一掷不计，换下一位。\textbf{叠骰、斜立：}全部在碗内但点数无法判定时，由同桌确认，6 颗全部重掷。

\vspace{7pt}\hrule height 0.4pt\vspace{5pt}
{\fontsize{10}{15}\selectfont
状元王：\underline{\hspace{2.4cm}}\quad 六颗骰面：\underline{\hspace{3.4cm}}\par}
\vspace{3pt}
{\fontsize{9}{12}\selectfont 本场以本页规则为准，桌内玩法见各桌规则纸。祝大家中秋快乐！\par}
\end{document}
